%%%%%%%%%%%%%%%%%%%%%%%%%%%%%%%%%%%%%%%%%%%%%%%%%%%%%%%%%%%%%%%%%%%%%%%%%%%
\chapter{Conclusion and Recommendations}\label{chap:recommendations}
%%%%%%%%%%%%%%%%%%%%%%%%%%%%%%%%%%%%%%%%%%%%%%%%%%%%%%%%%%%%%%%%%%%%%%%%%%%

\section{Conclusion}

The experiments presented in this thesis set out to determine whether the performance ceiling of a normal DQN agent in Stratego arises from insufficient training or from structural properties of the architecture itself. The evidence supports the latter interpretation. Across 37,213 episodes of curriculum training, the normal DQN exhibited a stable 22\% win rate, monotonically decreasing TD loss with no corresponding performance gain, Q-value overestimation, piece prediction accuracy frozen near 20\%, and a 38.5\% draw rate driven by passive piece shuffling. These observations are not independent failures. They form a coherent pattern: scalar value collapse prevents the agent from distinguishing aggressive high-variance moves from passive low-variance alternatives, exploration decay eliminates the mechanism for discovering new strategies, and the LSTM information bottleneck denies the agent the opponent model needed to evaluate combat outcomes.

The Rainbow DQN with AAREN targeted each of these bottlenecks with an architectural change. C51 distributional value estimation replaced scalar Q-learning. Noisy Networks replaced frozen $\epsilon$-greedy exploration with adaptive, state-dependent noise, which stayed active throughout training. AAREN replaced the LSTM history encoder; its piece-identity accuracy was 24.23\%, against 17.74\% for the LSTM. Multi-step returns shortened the effective bootstrap depth for sparse terminal rewards. In self-play the combined changes cut the draw rate from 50.04\% to 6.69\% and shortened the games. The self-play win rate rose from 24.97\% to 49.88\%, but in self-play both sides are the same agent, so that rise reflects the decisive share of games rather than greater strength (Section~\ref{finding:7}).

The transition from draw-dominated outcomes to decisive games is the most consequential behavioral shift observed in the data. The normal DQN converged on a policy that minimized short-term risk at the cost of long-term winning probability. The Rainbow agent stopped doing so in self-play. Against fixed heuristic opponents, however, the extended 75,047-episode run won 11.66\% of its games and drew 77.02\%. Of its wins, 99.96\% came by flag capture, and its average episode lasted 378 steps. Whether the agent learned to identify and pursue the flag, or only ended the games it won by that route, these figures cannot tell.

These results connect each architectural component to a failure mode of the baseline. They do not yet measure each component's contribution, for three reasons. The configurations differ in more than one component, and no ablation was run. Several of the measures used (the loss--win correlation, the logged Q-values, the self-play win rate) turned out not to measure what they were taken to measure (Chapter~\ref{chap:results}). The following sections present the individual components, the broader methodological contributions, and the remaining limitations.

\section{Recommended Architectural Enhancements}

\subsection{Value Estimation}

The Q-value overestimation and scalar value collapse identified in the training analysis directly motivated two Rainbow DQN enhancements formalized in the MARQ architecture. The Dueling Network decomposes action values into state-value and advantage streams. This design suits Stratego's action space, where many positional moves share similar strategic value and only a small subset of actions (attacks, flag-adjacent moves) carry distinct utility. The C51 distributional head provides multi-modal value estimation that differentiates risky decisive actions from safe passive alternatives. Scalar Q-values provably cannot make this distinction.

The empirical results are consistent with this design but do not isolate it. The correlation between loss and win rate shifted from 0.051 (normal DQN) to $-$0.787 (Rainbow DQN with AAREN). Both curves drift with training time, though, and in the later formal\_4 run the correlation shrank from $-0.49$ to $-0.14$ once that drift was removed, while the loss fell without any gain in evaluation (Section~\ref{finding:9}). The self-play draw rate dropped from 50.04\% to 6.69\%, but the 75k run against heuristics drew 77.02\%. Whether multi-modal value estimation by itself resolves the draw-seeking trap needs an ablation that changes only the value head.

\subsection{LSTM to AAREN Transition}

The LSTM prediction accuracy ceiling of approximately 18 to 22\% was the most direct evidence for the AAREN module's necessity. The LSTM-based history encoder has concrete advantages over a full DRQN architecture in replay compatibility, per-position granularity, and gradient isolation. The sequential information bottleneck inherent in any recurrent encoder, however, limits representational capacity over long game horizons. AAREN preserves all structural advantages of the separated encoder design while replacing the recurrent computation with attention-based history aggregation. This substitution addresses the compression bottleneck without modifying the replay buffer infrastructure or the DQN backbone.

AAREN achieved 24.23\% piece identity inference accuracy compared to 17.74\% for the LSTM, a 6.5 percentage point difference. This does not yet validate the transition. Always naming the most common piece type would score about 20 to 24\%, the two runs differed in more than the encoder, and the thesis code took AAREN's embedding at a padding step (Section~\ref{finding:11}). The separated encoder design itself is a useful structure. Independent hidden states for each of the 100 board positions produce a spatially structured $(64, 10, 10)$ embedding tensor that preserves piece-level temporal information. A monolithic DRQN hidden state would conflate action histories across all 40 opponent pieces into a single vector, losing spatial specificity. The per-position design also maintains full compatibility with standard experience replay by storing compact embedding snapshots alongside 15-channel board tensors, avoiding the sequence burn-in cost required by full DRQN architectures.

\subsection{Exploration Recovery}

The Noisy Net sigma remained at zero throughout the entire normal DQN training run. This confirmed that the $\epsilon$-greedy exploration schedule exhausted its utility without the agent discovering strategies beyond the initial local optimum. Noisy Networks parameterize exploration as learnable per-weight perturbations and provide state-dependent exploration that adapts to the agent's uncertainty rather than following a predetermined decay schedule.

The Rainbow agent's NoisyNet entropy rose steadily from 0.034 to 0.049 throughout training, so the learned exploration noise did not collapse. How much this exploration contributed to the change in outcomes is not known without an ablation.

\subsection{Long-Range Credit Assignment}

Single-step TD learning attenuates terminal rewards by $\gamma^{200} \approx 0.134$ for mid-game decisions. Five-step returns ($n = 5$), as specified in MARQ, shorten the effective bootstrap depth and push strategic outcomes back to earlier decision points faster. The steps-per-win figure (1,060 versus 2,697 for the normal DQN) does not test this. In self-play it follows the draw rate (Section~\ref{finding:10}), so whether multi-step returns accelerated reward propagation remains to be measured.


\section{Computer Science Contributions}

The following contributions apply to reinforcement learning systems beyond the Stratego domain. Stratego is the test domain, but the techniques are not domain-specific and apply to other partially observable or adversarial settings.

\subsection{Five-Phase Curriculum Design}

The curriculum transitions agents from full to partial observability in stages. Observability drops progressively while opponent difficulty rises and reward shaping decays. The intent is to mitigate the cold-start problem, so that agents do not need to learn basic mechanics and hidden-information reasoning at the same time. The later runs found a cost to this structure (Appendix~\ref{app:lessons}). Changing the observation and the shaping weight between phases made the network relearn the game at each switch. The current opponent schedule therefore keeps fog of war, the turn limit and the shaping fixed, and varies only the opponents.

\subsection{Embedding Snapshots for Replay Compatibility}

Hybrid architectures face a tension between replay accuracy and speed. Reconstructing full history sequences at replay time gives correct gradients but costs 4$\times$ throughput. Storing detached snapshot embeddings keeps throughput, but the DQN loss can then no longer train the encoder. The thesis code tried to restore that path with a ``gradient bridge'', adding $\epsilon \cdot W_{proj}(\mathbf{0})$ to the stored embeddings. That term equals the projection's bias, so it reached only that one parameter tensor and nothing that depends on the history (Chapter~\ref{chap:methodology}). The design that works is a separately supervised encoder: the snapshots are plain inputs to the Q-network, and the encoder is trained by its own loss on revealed pieces. Any encoder that is meant to learn from the value loss should be checked by listing which of its parameters actually receive a gradient.

\subsection{Expectamax Search Integration}

At test time, the system blends the DQN's positional Q-values with combat utilities derived from AAREN's rank-probability predictions through Expectamax search. This grounds neural value estimates in tactical reality and blocks overconfident attacks on hidden pieces. The approach transfers to any partially observable domain that combines learned value functions with probabilistic state inference.

\subsection{Potential-Based Reward Shaping with Phase Annealing}

PBRS guarantees policy invariance while providing dense signals during early training. A multiplicative schedule tapers the shaping coefficient to zero across curriculum phases, so the final training phase runs on terminal rewards alone. No manual cutoff is needed. The design avoids the reward-farming problem common in additive shaping.


\section{Limitations and Future Work}

The validation presented in this thesis was bounded by the training infrastructure, opponent diversity, and compute budget available during the research period. The following limitations qualify the conclusions drawn above and identify directions where additional work is needed.

The Rainbow DQN with AAREN produced far fewer draws than the normal DQN baseline in self-play. The two were never evaluated against each other or against the same fixed opponents, and against heuristic opponents the Rainbow agent won 11.66\% of its games. Performance against opponents beyond the current heuristic and self-play pool may reveal gaps in strategic reasoning that the existing evaluation does not expose.

Combining C51, Noisy Nets, Dueling Networks, AAREN, and $n$-step returns in a multi-agent self-play loop introduces interaction effects that single-component theory does not predict. The extended 75k training run remained in the partial observability phase against heuristic opponents for its entire duration, suggesting that curriculum promotion thresholds require further tuning for long-horizon training stability.

The full MARQ pipeline is heavier than the normal DQN. Each checkpoint occupies approximately 433\,MB compared to 39\,MB, an 11-fold increase. While sample efficiency improved (5.17 million steps versus 10.31 million for the normal DQN), wall-clock time per episode is higher, reducing total episode count under a fixed compute budget.

Ablation studies are needed. The data does not yet show that the combined architecture plays better than the normal DQN, nor what each component contributes. Future work should test each component in isolation, evaluated in greedy games against the same fixed opponents, with enough games for a difference to exceed its noise. AAREN in particular needs paired runs with and without it, and a majority-class baseline for its accuracy. Beyond fixed opponents, the agent's exploitability should be measured by training a fresh exploiter against the frozen policy. A deterministic greedy policy can be exploited in a game of hidden information, and methods built for such games (NFSP, PSRO, R-NaD) add an averaged or regularized policy for that reason.
